\SourceLeaf{198}
\Needspace{8\baselineskip}
\pdfbookmark[1]{Seeber, positive ternäre quadratische Formen. Besprechung vom 9. Juli 1831}{notice-1831-07-09}
\begin{center}
\rule{.86\linewidth}{.3pt}\par
Göttingische gelehrte Anzeigen.\quad 9. Juli 1831.\par
\rule{.86\linewidth}{.3pt}
\end{center}
\emph{Untersuchungen über die Eigenschaften der positiven ternären quadratischen Formen von} \textsc{Ludwig August Seeber}, \emph{Dr. der Philosophie, ordentlicher Professor der Physik an der Universität in Freiburg. Freiburg im Breisgau} 1831. (248 S. in 4.)

Die Funktionen zweier unbestimmter Größen $x$ und $y$ von der Gestalt $axx+2bxy+cyy$, wobei $a,b,c$ bestimmte ganze Zahlen bezeichnen, bilden bekanntlich unter dem Namen \emph{quadratische Formen} oder, wenn eine weitere Unterscheidung nötig ist, \emph{binäre quadratische Formen} einen der interessantesten und reichhaltigsten Gegenstände der höheren Arithmetik. Dabei stellen sich zunächst folgende Aufgaben: zu entscheiden, ob eine solche gegebene Form eine andere $a'x'x'+2b'x'y'+c'y'y'$ enthält, also durch eine Substitution $x=\alpha x'+\beta y'$, $y=\gamma x'+\delta y'$, in der $\alpha,\beta,\gamma,\delta$ ganze Zahlen sind, in diese überführt werden kann; ob eine solche Beziehung zwischen zwei Formen gegenseitig ist, in welchem Fall die Formen äquivalent heißen; ferner in beiden Fällen alle möglichen Umformungen der einen in die andere anzugeben; schließlich alle möglichen Darstellungen einer gegebenen ganzen Zahl durch eine gegebene Form mit ganzzahligen Werten der unbestimmten Größen zu finden. Diese Aufgaben sind in den \emph{Disquisitiones Arithmeticae} vollständig gelöst, bilden aber nur den bei Weitem kleineren Teil des dortigen Abschnitts über quadratische Formen. Die folgenden feineren Untersuchungen erfordern\textsuperscript{.} zum Teil eine vorläufige Bearbeitung eines um eine Stufe höheren und viel größere Schwierigkeiten bietenden Gebiets: der Lehre von ähnlichen Funktionen dreier unbestimmter Größen $x,y,z$, die also die Gestalt $axx+byy+czz+2a'yz+2b'xz+2c'xy$ haben und ternäre quadratische Formen heißen. Die Lösung der Hauptaufgaben zu diesen ternären Formen ist in dem genannten Werk entwickelt, jedoch nur so weit, wie es für den angegebenen Zweck nötig war. Nach einem Abstand von dreißig Jahren hat nun der Verfasser des vorliegenden Werks diese Untersuchungen als Erster wieder aufgenommen und für eine Hauptgattung der ternären Formen, nämlich die positiven, das in den \emph{Disquisitiones }
\SourceLeaf{199}
\emph{Arithmeticae} unvollendet Gelassene vervollständigt. Denjenigen, die sich eingehender mit der höheren Arithmetik beschäftigt haben, könnten wir mit wenigen Worten sagen, was in dem vorliegenden Werk neu geleistet wird. Um aber auch anderen verständlich zu sein, müssen wir uns etwas mehr Ausführlichkeit gestatten. Das tun wir umso lieber, als diese Untersuchungen auch außerhalb der höheren Arithmetik ein eigenes Interesse haben.

Die Eigenschaften einer binären Form $axx+2bxy+cyy$ hängen vor allem von der Zahl $bb-ac$ ab, die daher die Determinante dieser Form heißt. Zwei äquivalente Formen haben immer gleiche Determinanten. Aber nicht alle Formen mit einer gegebenen Determinante sind deshalb schon äquivalent. Vielmehr zerfallen solche Formen in eine kleinere oder größere, stets jedoch endliche Anzahl von Klassen, sodass Formen derselben Klasse untereinander äquivalent sind, Formen verschiedener Klassen dagegen nicht. Durch Formen mit positiver Determinante lassen sich positive und negative Zahlen gleichermaßen darstellen; durch Formen mit negativer Determinante dagegen nur Zahlen, die dasselbe Vorzeichen wie $a$ und $c$ haben. Daher unterscheidet man hier positive und negative Formen. Die einfachsten Formen in jeder Klasse haben bestimmte Kriterien, heißen reduzierte Formen und können als Vertreter der ganzen Klasse betrachtet werden.

Ähnliche Verhältnisse für die ternären Formen sind in den \emph{Disquisitiones Arithmeticae} nachgewiesen. Die Determinante der ternären Form
\[
axx+byy+czz+2a'yz+2b'xz+2c'xy
\]
ist die Zahl
\[
aa'a'+bb'b'+cc'c'-abc-2a'b'c'
\]
Auch hier ist die Gleichheit der Determinanten für die Äquivalenz zweier Formen notwendig, aber nicht hinreichend. Vielmehr zerfallen sämtliche Formen mit einer bestimmten Determinante in eine endliche Anzahl von Klassen, in denen jeweils die einfachsten Formen reduziert heißen und gleichsam alle übrigen vertreten können. Mit dem Unterschied zwischen positiven und negativen Formen verhält es sich hier aber anders als bei den binären Formen. Für jede gegebene Determinante, ob positiv oder negativ, gibt es einerseits Formen, durch die
\SourceLeaf{200}
positive und negative Zahlen gleichermaßen darstellbar sind (indifferente Formen), andererseits Formen, durch die sich entweder nur positive oder nur negative Zahlen darstellen lassen (positive oder negative Formen). Positive Formen gibt es aber nur für negative Determinanten und negative Formen nur für positive. Im Übrigen ist von selbst klar, dass die Eigenschaft einer Form, indifferent, positiv oder negativ zu sein, zugleich der ganzen Klasse zukommt, zu der sie gehört. Das vorliegende Werk beschränkt sich auf positive Formen, deren Determinanten also negativ sein müssen. Alles, was für diese gilt, überträgt sich offensichtlich von selbst auf die negativen Formen, während die im Werk ganz ausgeschlossenen indifferenten Formen eine völlig andere Behandlung erfordern.

In den \emph{Disquisitiones Arithmeticae} war, wie schon erwähnt, die Theorie der ternären Formen nur so weit entwickelt, wie es für den dortigen Zweck nötig war. Daher war die Aufgabe, die Äquivalenz zweier gegebener ternärer Formen zu entscheiden, noch nicht in voller Allgemeinheit gelöst. Zwar war dort gezeigt worden, wie man zu jeder vorgegebenen Form eine äquivalente Form der einfachsten Art findet und dass es für jede gegebene Determinante nur endlich viele solche reduzierten Formen geben kann. Da aber jede Klasse mehrere solche reduzierten Formen enthält, die sich nicht in allen Fällen \emph{sofort} als äquivalent erweisen, fehlte noch ein Kriterium, mit dem sich ihre Äquivalenz oder Nichtäquivalenz sicher erkennen lässt. Diesem Bedürfnis hat der Verfasser des vorliegenden Werks für die positiven Formen nun vollständig und mit vorbildlicher Gründlichkeit entsprochen. Sein Verfahren ist allerdings etwas anders eingekleidet, als wir die Sache gerade dargestellt haben und wie sie sich darstellen müsste, wenn man in den Begriff der reduzierten positiven Formen nur die wesentlichsten Bedingungen größter Einfachheit aufnimmt. Bei positiven Formen bestehen diese darin, dass die ihrer Natur nach positiven Zahlen $a,b,c$ nicht kleiner sein dürfen als jeweils $b'$ oder $c'$, $a'$ oder $c'$, $a'$ oder $b'$, ohne Rücksicht auf die Vorzeichen. Herr \textsc{Seeber} hat den Begriff der reduzierten Formen um solche Bedingungen ergänzt, dass es in jeder Klasse immer nur eine solche Form geben kann, aber eine geben muss. Wegen eines schönen, von Herrn \textsc{Seeber} aus beobachteten Fällen gewonnenen Satzes, der weiter unten noch erwähnt wird, geben wir hier die Hauptbedingungen an, die Herr S. in den Begriff der reduzierten Formen
\SourceLeaf{201}
aufgenommen hat: 1) Unter den Zahlen $a',b',c'$ dürfen nicht zwei entgegengesetzte Vorzeichen haben; 2) ohne Rücksicht auf das Vorzeichen dürfen $2b'$ und $2c'$ nicht größer als $a$ sein, ferner $a$ und $2a'$ nicht größer als $b$ und $b$ nicht größer als $c$; 3) wenn $a',b',c'$ zugleich negativ sind, darf die doppelte Summe dieser Zahlen nicht größer als $a+b$ sein. Die übrigen für einige Sonderfälle hinzukommenden Bedingungen können wir hier übergehen.

Den Hauptinhalt des Werks bilden nun zunächst die Lösung der Aufgabe, zu jeder gegebenen positiven Form eine äquivalente zu finden, die nach der festgelegten Definition den Charakter einer reduzierten Form hat, und dann der strenge Beweis des Satzes, dass zwei nicht identische reduzierte Formen nicht äquivalent sein können, oder, was dasselbe ist, dass jede Klasse nur eine reduzierte Form enthält. Dem Geist der Gründlichkeit, mit dem diese Gegenstände durchgeführt sind, müssen wir volle Gerechtigkeit widerfahren lassen. Und wenn wir dabei die sehr große, vielleicht manchen abschreckende Weitläufigkeit bedauern müssen, da die Lösung des Problems 41 Seiten und der Beweis des Satzes 91 Seiten einnimmt, so soll dies doch keineswegs als Tadel verstanden werden. Wenn ein schwieriges Problem zu lösen oder ein schwieriger Satz zu beweisen ist, besteht der erste und mit gebührendem Dank anzuerkennende Schritt stets darin, dass überhaupt eine Lösung oder ein Beweis gefunden wird. Die Frage, ob dies nicht leichter und einfacher hätte geschehen können, bleibt so lange müßig, wie die Möglichkeit nicht zugleich durch die Ausführung entschieden wird. Wir halten es daher für verfrüht, hier bei dieser Frage zu verweilen. --- Der übrige Teil des Werks enthält hauptsächlich die mit gleicher Gründlichkeit ausgeführten Lösungen folgender Aufgaben: zu entscheiden, ob eine gegebene Form eine andere gegebene, ihr nicht äquivalente Form enthält; alle möglichen Transformationen einer gegebenen Form in eine gegebene äquivalente oder nur in ihr enthaltene Form zu finden; schließlich für eine gegebene Determinante alle möglichen Klassen positiver ternärer Formen anzugeben.

Wir müssen noch bemerken, dass Herr \textsc{Seeber} die Gestalt der ternären Formen etwas anders gefasst hat als die \emph{Disquisitiones Arithmeticae}, wo die Koeffizienten der Produkte $yz$, $xz$, $xy$ mit Bedacht als gerade Zahlen vorausgesetzt waren. Herr S. lässt dagegen auch ungerade zu und bezeichnet daher
\SourceLeaf{202}
mit $a',b',c'$, was oben mit $2a'$, $2b'$, $2c'$ bezeichnet war. Die dadurch erreichte größere Allgemeinheit ist offensichtlich nur scheinbar oder doch überflüssig, da sich alles, was über solche Formen mit ungeraden Koeffizienten gesagt werden kann, auch von selbst ergibt, wenn man stattdessen ihr Doppeltes betrachtet. Wir können diese Änderung daher nicht billigen, zumal sie auch einen gewissen Verlust an Einfachheit mit sich bringt. Eine Folge davon ist, dass das, was Herr \textsc{Seeber} Determinante nennt, immer das Vierfache der Zahl ist, die in den \emph{Disquisitiones Arithmeticae} diesen Namen trägt. In dieser Besprechung haben wir die Terminologie der \emph{Disquisitiones Arithmeticae} beibehalten.

Beim zuletzt genannten Problem, für jede gegebene Determinante alle möglichen reduzierten Formen anzugeben, hat Herr \textsc{Seeber} einen Satz benutzt, um Grenzen für die ersten drei Koeffizienten zu erhalten. Danach kann ihr Produkt $abc$ nicht größer als das Dreifache der Determinante sein. Diesen Satz hat Herr \textsc{Seeber} streng bewiesen. In der Vorrede bemerkt er jedoch, er habe unter mehr als 600 untersuchten Fällen keinen einzigen gefunden, in dem jenes Produkt das Doppelte der Determinante überschritten hätte. Daher hält er es für höchst wahrscheinlich, dass diese engere Grenze allgemein gilt; einen strengen Beweis habe er aber nicht finden können. Da dieser von Herrn \textsc{Seeber} aus beobachteten Fällen gewonnene Satz sowohl an sich bemerkenswert als auch für eine kürzere Lösung der genannten Aufgabe wichtig ist, wollen wir hier einen sehr einfachen Beweis hinzufügen und damit auch unsererseits in dieser Besprechung zur Vervollkommnung der Theorie beitragen. Dabei sind zwei Fälle zu unterscheiden.

\Needspace{22\baselineskip}\begingroup\postdisplaypenalty=10000 I. Wenn keine der Zahlen $a',b',c'$ negativ ist, setze man
\[
\begin{aligned}
b-2a'&=d,&c-2b'&=e,&a-2c'&=f\\
c-2a'&=g,&a-2b'&=h,&b-2c'&=i
\end{aligned}
\]
Aus der Definition der reduzierten positiven Formen folgt sofort: Wenn
\[
axx+byy+czz+2a'yz+2b'xz+2c'xy
\]
eine solche Form ist, ist keine dieser sechs Zahlen negativ. Ebenso versteht sich von selbst, dass $a,b,c$ positiv sind. Bezeichnet man nun die (negative)
\SourceLeaf{203}
Determinante der Form mit $-D$, so gilt die identische Gleichung, von der man sich leicht durch Ausmultiplizieren überzeugt,
\[
2D-abc=aa'd+bb'e+cc'f+a'hi+b'gi+c'gh+ghi
\]
in der keiner der sieben Terme rechts negativ sein kann. Folglich ist $abc$ nicht größer als $2D$. Dasselbe folgt auf die gleiche Weise aus der identischen Gleichung
\[
2D-abc=aa'g+bb'h+cc'i+a'ef+b'df+c'de+def
\]\endgroup

\Needspace{23\baselineskip}II. Wenn keine der Zahlen $a',b',c'$ positiv ist, setze man
\[
\begin{aligned}
b+2a'&=d,&c+2b'&=e,&a+2c'&=f\\
c+2a'&=g,&a+2b'&=h,&b+2c'&=i
\end{aligned}
\]
\[
\begin{aligned}
b+c+2a'+2b'+2c'&=k\\
a+c+2a'+2b'+2c'&=l\\
a+b+2a'+2b'+2c'&=m
\end{aligned}
\]
und die Determinante der Form wie zuvor $=-D$. Nach der Definition der reduzierten positiven Formen kann keine der neun Zahlen $d,\allowbreak e,\allowbreak f,\allowbreak g,\allowbreak h,\allowbreak i,\allowbreak k,\allowbreak l,\allowbreak m$ negativ sein. Aus der identischen Gleichung
\[
\begin{aligned}
6D-3abc={}&-aa'(d+2k)-bb'(e+2l)-cc'(f+2m)\\
&-a'hi-b'gi-c'gh+def+2ghi
\end{aligned}
\]
in der alle Terme rechts positiv oder null sind, weil $a',b',c'$ nicht positiv, sondern negativ oder null sind, ergibt sich daher, dass $3abc$ nicht größer als $6D$ oder $abc$ nicht größer als $2D$ sein kann. Dasselbe folgt ebenso aus der identischen Gleichung
\[
\begin{aligned}
6D-3abc={}&-aa'(g+2k)-bb'(h+2l)-cc'(i+2m)\\
&-a'ef-b'df-c'de+2def+ghi
\end{aligned}
\]
Beide Gleichungen sind symmetrisch. Verzichtet man auf vollständige Symmetrie, so lässt sich der Beweis mit noch weniger Termen führen, zum Beispiel durch die identische Gleichung
\[
\begin{aligned}
8D-4abc={}&-2aa'(g+k)-2bb'(e+l)-4cc'm\\
&+(c+e)df+(c+g)hi
\end{aligned}
\]

\SourceLeaf{204}
Wir wollen nun noch einiges über die Bedeutung der positiven binären und ternären quadratischen Formen außerhalb der höheren Arithmetik hinzufügen. Die negativen Formen gesondert zu behandeln ist unnötig, und die indifferenten entziehen sich dieser Behandlung ganz.

\Needspace{9\baselineskip}Die positive binäre Form $axx+2bxy+cyy$ stellt allgemein das Quadrat der Entfernung zweier unbestimmter Punkte in einer Ebene dar, deren Koordinaten sich in Bezug auf zwei gegeneinander geneigte Achsen, deren Winkel den Kosinus $=\frac{b}{\surd ac}$ hat, um $x\surd a$, $y\surd c$ unterscheiden. Wenn $x$ und $y$ also nur \emph{ganze} Zahlen bedeuten sollen, bezieht sich die Form auf ein System von in Parallelogrammen angeordneten Punkten, die in den Schnittpunkten zweier Systeme paralleler Geraden liegen. Die Geraden jedes Systems haben gleiche Abstände voneinander: Die Abstände des einen Systems, parallel zu den Geraden des zweiten gemessen, sind $=\surd a$, die Abstände des anderen, parallel zu den Geraden des ersten gemessen, $=\surd c$; die Neigung der beiden Systeme ist die oben angegebene. So erscheint die Ebene vollständig in gleiche Parallelogramme geteilt, deren Eckpunkte das Punktesystem bilden, ohne dass irgendeiner dieser Punkte ins Innere eines Parallelogramms fallen kann. Die mit positivem Vorzeichen genommene Determinante, also $ac-bb$, bedeutet das Quadrat des Flächeninhalts eines elementaren Parallelogramms. Ein und dasselbe System solcher Punkte kann auf unendlich viele verschiedene Arten in Parallelogramme eingeteilt und daher auf ebenso viele verschiedene Formen bezogen werden. Alle diese verschiedenen Formen sind aber, wie es in der Fachsprache heißt, äquivalent, und der Flächeninhalt eines elementaren Parallelogramms bleibt immer derselbe. Zwei nicht äquivalente Formen, von denen aber die eine die andere enthält, beziehen sich auf dasselbe Punktesystem: die erste Form auf das ganze System, die zweite auf einen Teil. Zwei Formen, die in der Fachsprache uneigentlich äquivalent (improprie aequivalentes) heißen, beziehen sich auf zwei gleiche, aber entgegengesetzt liegende Punktesysteme, wobei man sich die Ebene umgekehrt gelegt denkt, und so weiter.

Auf die gleiche Weise bedeutet die positive ternäre Form
\[
axx+byy+czz+2a'yz+2b'xz+2c'xy
\]
allgemein das Quadrat der Entfernung zweier unbestimmter Punkte im Raum, deren Koordinaten in Bezug auf drei Achsen (1), (2), (3) die Unterschiede $x\surd a$, $y\surd b$, $z\surd c$
\SourceLeaf{205}
ergeben: Die Kosinus der Winkel zwischen den Achsen (2) und (3), (1) und (3), (1) und (2) sind jeweils $\frac{a'}{\surd bc}$, $\frac{b'}{\surd ac}$, $\frac{c'}{\surd ab}$. Wenn $x,y,z$ hier nur ganze Zahlen bedeuten sollen, bezieht sich die Form auf ein System von in Parallelepipeden angeordneten Punkten, die also durch die Schnittpunkte dreier Systeme paralleler, gleichabständiger Ebenen entstehen. Der ganze Raum erscheint so vollständig in gleiche Parallelepipede geteilt, deren Eckpunkte jenes Punktesystem bilden. Das Quadrat des Rauminhalts eines elementaren Parallelepipeds ist gleich der mit positivem Vorzeichen genommenen\textsuperscript{.} Determinante der ternären Form. Äquivalente Formen vertreten ein und dasselbe Punktesystem, nur auf andere Achsen oder Grundebenen bezogen. Ebenso finden alle anderen Hauptbestandteile der Theorie der ternären Formen hier ihre geometrische Bedeutung: das Enthaltensein einer Form in einer anderen, die Darstellung einer bestimmten Zahl oder einer unbestimmten binären Form durch eine ternäre, die Lehre von den zugeordneten ternären Formen (formae adiunctae), der Wegfall des Unterschieds zwischen eigentlicher und uneigentlicher Äquivalenz, das Wesen der reduzierten Formen und so weiter. Wir müssen uns jedoch auf die obigen Andeutungen beschränken, zumal das vorliegende Werk, das die ternären Formen ausschließlich aus rein arithmetischer Sicht betrachtet, nur mittelbar dazu Anlass gegeben hat. Man wird daraus wenigstens erkennen, welch reiches Feld sich hier für Untersuchungen eröffnet, die nicht nur für sich ein hohes theoretisches Interesse haben, sondern auch zu einer ebenso bequemen wie allgemeinen Behandlung aller Beziehungen zwischen den Kristallformen dienen können. Auf die Einzelheiten dieser Anwendung einzugehen ist hier nicht der Ort. Wir dürfen jedoch nicht übergehen, dass, obwohl ursprünglich angenommen wurde, $a,b,c,a',b',c'$ bezeichneten ganze Zahlen, der größte Teil der Lehre von den ternären Formen und insbesondere das für jene Anwendung Erforderliche auch unabhängig von dieser Voraussetzung gültig bleibt. Tatsächlich führen zwar die Angaben von \textsc{Haüy} bei den meisten Kristallarten auf sehr einfache ganzzahlige Werte der Koeffizienten der ternären Formen, die sich auf die entsprechende Anordnung des Punktesystems beziehen. Die genaueren späteren Messungen von \textsc{Wollaston}, \textsc{Malus}, \textsc{Biot}, \textsc{Kupffer} und anderen widersprechen dem jedoch und lassen es zweifelhaft erscheinen, ob rationale Verhältnisse dieser Koeffizienten überall der Natur entsprechen. Wenn man in der Theorie aber die
\SourceLeaf{206}
Beschränkung auf ganzzahlige Koeffizienten nicht weglassen will, lassen sich jedenfalls immer ganze Zahlen finden, die den Messergebnissen beliebig nahe kommen, da es dabei nur auf die Verhältnisse untereinander und nicht auf absolute Werte ankommt.

Schließlich wollen wir dem oben angeführten Satz von \textsc{Seeber} noch seine geometrische Bedeutung geben. Wenn ein Parallelepiped so beschaffen ist, dass keine seiner zwölf Kanten, von denen jeweils vier gleich sind, größer ist als irgendeine der zwölf Diagonalen der Seitenflächen, die paarweise gleich sind, oder irgendeine der vier Diagonalen des Parallelepipeds, dann ist sein mit $\surd2$ multiplizierter Rauminhalt nicht kleiner als der Rauminhalt eines aus denselben Kanten gebildeten rechtwinkligen Parallelepipeds.
\BookRule{.28\linewidth}
